\centerline{\bf \Huge Турбо тренировка V}
\centerline{\bf \Huge }

\begin{flushright}
        \scriptsize{всё плохо...}
\end{flushright}

\problem{1}
Анджур Аркадьевич так увлёкся задачами на разрезания, что решил открыть свой клуб "Разрезателей". Каждому участнику клуба он выдал по прямоугольнику (всем одного и того же размера) и попросил их разрезать его на одинаковые квадратики. Оказалось, что общее число квадратиков --- простое число. Докажите, что исходные прямоугольники тоже были квадратами.

\problem{2}
На полу нарисован правильный 99-угольник. Эльзята и Олежа играют в игру, делая ходы по очереди. Первой ходит Эльзята; своим первым ходом она кладёт в одну из вершин лампочку, причём она может быть как включённой так и выключенной (Эльзята сама выбирает в каком она будет состоянии). Затем каждый из игроков своим ходом может положить включенную или выключенную лампочку в любую пустую вершину, соседнюю с той вершиной, где лампочка уже стоит. Олежа выигрывает, если после того как все вершины будут заняты лампочками (то есть после 99 ходов) найдётся равносторонний треугольник (его вершины это вершины многоугольника) у которого либо все вершины горят, либо все вершины не горят. Может ли Эльзята помешать Олеже?

\problem{3}
Пусть $m, n \in \mathbb{N}.$ Рассмотрим всевозможные попарные произведения чисел $m, n, m+2, n+2$. Какое наибольшее количество квадратов может быть среди этих 6 чисел?

\problem{4}
Дан остроугольный треугольник $ABC$ с углом $B$ равным $60^{\circ}$. Внутри отмечена точка $X$ такая, что $\angle AXB =\angle BXC=\angle AXC = 120^{\circ}$. Пусть $M$ --- центр масс треугольника. Прямая $XM$ пересекает описанную окружность треугольника $AXC$ в точке $Y$. Чему равно $\dfrac{XM}{MY}$